\begin{helper}[Arbitrary marked-set injection image tail]
Fix natural numbers \(n,m,q\), a marked set
\(K\subseteq \operatorname{Fin}(m)\times\operatorname{Fin}(m)\) with
\(|K|=q\), and \(0<\delta\le 1/2\).  Assume \(n\le m^2\), and put
\[
  \mu = n\frac{q}{m^2}.
\]
Then the number of injections
\(\phi:\operatorname{Fin}(n)\hookrightarrow
\operatorname{Fin}(m)\times\operatorname{Fin}(m)\) whose image hits \(K\)
outside the interval \([(1-\delta)\mu,(1+\delta)\mu]\) is at most
\[
  \left(
    \exp\!\left(-(\delta+(1-\delta)\log(1-\delta))\mu\right)
    +
    \exp\!\left(-(((1+\delta)\log(1+\delta)-\delta)\mu)\right)
  \right)
  \left|\operatorname{Inj}\!\left(
    \operatorname{Fin}(n),
    \operatorname{Fin}(m)\times\operatorname{Fin}(m)
  \right)\right|.
\]
\end{helper}
\begin{proof}
Transport the product population
\(\operatorname{Fin}(m)\times\operatorname{Fin}(m)\) to
\(\operatorname{Fin}(m^2)\) using the standard product equivalence.  The
transported marked set has cardinality \(q\), and \(q\le m^2\) because it is
the cardinality of a subset of the product population.  For a uniformly chosen
\(n\)-subset of \(\operatorname{Fin}(m^2)\), the upper Chernoff bound
\(\noderef{HypergeometricCountingUpperChernoffBound}\) gives
\[
  \Pr\{|S\cap K|>(1+\delta)\mu\}
  \le
  \exp\!\left(-(((1+\delta)\log(1+\delta)-\delta)\mu)\right),
\]
and the lower-tail bound
\(\noderef{HypergeometricCountingLowerTailBound}\) gives
\[
  \Pr\{|S\cap K|<(1-\delta)\mu\}
  \le
  \exp\!\left(-(\delta+(1-\delta)\log(1-\delta))\mu\right).
\]
The nodes \(\noderef{FinProductTailCountEquiv}\) and
\(\noderef{FinProductLowerTailCountEquiv}\) identify these transported
subset counts with the corresponding product-population subset counts.
Finally, \(\noderef{UniformInjectionImage}\) says that every \(n\)-subset of
the product population is hit by the same number of injections, so multiplying
the two subset probabilities by the total number of injections gives the
corresponding upper and lower injection-count bounds.  The injections outside
the displayed interval are contained in the union of the upper-tail and
lower-tail injection events, so the finite union bound gives the displayed
two-sided estimate.
\end{proof}
